\documentclass[11pt]{article}

\usepackage[margin=1in]{geometry}
\usepackage{amsmath, amssymb, amsthm}
\usepackage{mathtools}
\usepackage{natbib}
\usepackage{hyperref}
\usepackage{xcolor}

\hypersetup{
  colorlinks=true,
  linkcolor=black,
  citecolor=blue!50!black,
  urlcolor=blue!50!black
}

\title{{{title}}}
\author{Literature review entry}
\date{Studied: {{date}} \\ \texttt{{{arxiv_id}}}}

\begin{document}
\maketitle

\begin{abstract}
{{abstract_paragraph}}
\end{abstract}

\section{Background and Motivation}
{{background_paragraph}}

\section{Technical Approach}
\label{sec:method}

{{method_intro_paragraph}}

\paragraph{Definitions.}
{{definitions_paragraph_with_inline_math}}

\paragraph{Central object.}
The central quantity studied is
\begin{equation}
{{central_equation}}
\label{eq:central}
\end{equation}
where {{symbol_glossary}}.

\paragraph{Method.}
{{method_walkthrough}}

\section{Key Results}
{{results_paragraph}}

\section{Position in the Literature}
{{position_paragraph_with_citations}}

\section{Limitations and Open Questions}
{{limitations_paragraph}}

\section*{Sandbox Probe}
{{sandbox_paragraph_optional}}

\bibliographystyle{plainnat}
\bibliography{references}

\end{document}
